\documentclass[12pt, a4paper]{book}

\usepackage[T1]{fontenc}
\usepackage[utf8]{inputenc}
\usepackage[charter, cal=cmcal]{mathdesign}
\usepackage[scaled=0.92]{helvet}
\usepackage{geometry}
\geometry{a4paper, top=30mm, bottom=30mm, inner=32mm, outer=25mm, marginparwidth=18mm}

\usepackage{amsmath, amsthm}
\usepackage[version=4]{mhchem}
\usepackage{graphicx}
\usepackage{xcolor}
\usepackage{tikz}
\usetikzlibrary{arrows.meta, shapes.geometric, positioning, calc, backgrounds}
\usepackage{pgfplots}
\pgfplotsset{compat=1.18}
\usepackage{tcolorbox}
\tcbuselibrary{skins, breakable, theorems}
\usepackage{titlesec}
\usepackage{fancyhdr}
\usepackage{enumitem}
\usepackage{booktabs}
\usepackage{array}
\usepackage{parskip}
\usepackage{siunitx}
\usepackage[colorlinks=true, linkcolor=BrickRed,
            urlcolor=NavyBlue, citecolor=ForestGreen]{hyperref}

% ── Colours (Charcoal / Royal Blue) ──────────────────────────────────────────
\definecolor{JourneyColor}{HTML}{0F172A}
\definecolor{ModuleColor}{HTML}{2563EB}
\definecolor{PuzzleColor}{HTML}{D97706}
\definecolor{DiscoveryColor}{HTML}{3B82F6}
\definecolor{WarningColor}{HTML}{991B1B}
\definecolor{ProblemColor}{HTML}{6B21A8}
\definecolor{ReflectColor}{HTML}{0D9488}
\definecolor{HistoryColor}{HTML}{78350F}

% ── Boxes ─────────────────────────────────────────────────────────────────────
\newtcolorbox{puzzlebox}[1][]{enhanced,breakable,
  colback=PuzzleColor!6,colframe=PuzzleColor,fonttitle=\bfseries\sffamily,
  title={The Puzzle},attach boxed title to top left={yshift=-2mm,xshift=6mm},
  boxed title style={colback=PuzzleColor},borderline west={3pt}{0pt}{PuzzleColor},#1}

\newtcolorbox{discoverybox}[1][]{enhanced,breakable,
  colback=DiscoveryColor!8,colframe=DiscoveryColor,fonttitle=\bfseries\sffamily,
  title={Discovery},attach boxed title to top left={yshift=-2mm,xshift=6mm},
  boxed title style={colback=DiscoveryColor},borderline west={3pt}{0pt}{DiscoveryColor},#1}

\newtcolorbox{problembox}[1][]{enhanced,breakable,
  colback=ProblemColor!5,colframe=ProblemColor,fonttitle=\bfseries\sffamily,
  title={Problems},attach boxed title to top left={yshift=-2mm,xshift=6mm},
  boxed title style={colback=ProblemColor},borderline west={3pt}{0pt}{ProblemColor},#1}

\newtcolorbox{deepbox}[1][]{enhanced,breakable,
  colback=JourneyColor!5,colframe=JourneyColor,fonttitle=\bfseries\sffamily,
  title={Deep Problems},attach boxed title to top left={yshift=-2mm,xshift=6mm},
  boxed title style={colback=JourneyColor},borderline west={4pt}{0pt}{JourneyColor},#1}

\newtcolorbox{historybox}[1][]{enhanced,breakable,
  colback=HistoryColor!7,colframe=HistoryColor!60,fonttitle=\bfseries\sffamily\small,
  title={History Note},attach boxed title to top left={yshift=-2mm,xshift=6mm},
  boxed title style={colback=HistoryColor!80},#1}

\newtcolorbox{reflectionbox}[1][]{enhanced,breakable,
  colback=ReflectColor!6,colframe=ReflectColor,fonttitle=\bfseries\sffamily,
  title={Reflection},attach boxed title to top left={yshift=-2mm,xshift=6mm},
  boxed title style={colback=ReflectColor},borderline west={3pt}{0pt}{ReflectColor},#1}

\newtcolorbox{namedconcept}[2][]{enhanced,
  colback=ModuleColor!6,colframe=ModuleColor,fonttitle=\bfseries\sffamily\large,
  title={#2},attach boxed title to top left={yshift=-2mm,xshift=6mm},
  boxed title style={colback=ModuleColor},#1}

% ── Section formatting ────────────────────────────────────────────────────────
\titleformat{\chapter}[display]
  {\normalfont\sffamily\huge\bfseries\color{JourneyColor}}
  {\textcolor{JourneyColor!60}{\chaptertitlename\ \thechapter}}{10pt}{\Huge}
  [\vspace{2pt}\textcolor{JourneyColor}{\titlerule[1.5pt]}]
\titleformat{\section}
  {\normalfont\sffamily\Large\bfseries\color{ModuleColor}}
  {\thesection}{1em}{}[\textcolor{ModuleColor!40}{\titlerule[0.6pt]}]
\titleformat{\subsection}
  {\normalfont\sffamily\large\bfseries\color{JourneyColor!80}}{\thesubsection}{1em}{}
\titleformat{\subsubsection}
  {\normalfont\sffamily\normalsize\bfseries}{\thesubsubsection}{1em}{}

\pagestyle{fancy}
\fancyhf{}
\fancyhead[LE]{\small\sffamily\textcolor{JourneyColor}{\leftmark}}
\fancyhead[RO]{\small\sffamily\textcolor{ModuleColor}{\rightmark}}
\fancyfoot[C]{\small\thepage}
\renewcommand{\headrulewidth}{0.4pt}

\newcommand{\newterm}[1]{\textit{\textcolor{DiscoveryColor}{#1}}}
\newcommand{\important}[1]{\textcolor{ModuleColor}{\textbf{#1}}}

\begin{document}

\setcounter{chapter}{5}
\chapter[Journey 6 --- Distributed Systems and Concurrency]
        {Journey Six\\[8pt]
         \textcolor{ModuleColor}{\large Distributed Systems and Concurrency}}
\setcounter{section}{0}

\section{Module 6.1 --- Concurrency, Mutual Exclusion, and Locks}
\label{sec:mod6.1}

\subsection*{Motivation}

We completed in Journey 5 the analysis and optimization of algorithms, proving bounds on search trees and network flows.
All those models assumed a single thread of execution running sequentially.
But modern processors are composed of multiple cores executing concurrently.
We must study the physics of cooperation.

When multiple threads access shared memory concurrently, execution steps interleave.
This can cause race conditions, resulting in corrupt data.
This module formalizes how threads coordinate access to shared state.
We will study critical sections and mutual exclusion, formulate the four Coffman conditions for deadlocks, and explore lock-free concurrency using atomic hardware primitives like Compare-And-Swap (CAS).

\subsection{The Shared Counter}

We analyze how interleaved execution causes race conditions.

\begin{puzzlebox}
\textbf{The Interleaved Increment.}

Suppose we have a shared variable in memory: `counter = 0`.
We spawn two threads that run concurrently on two CPU cores. Both threads execute the same code block $10^6$ times:
`counter = counter + 1`
When both threads finish, we expect the value of `counter` to be exactly $2 \times 10^6$.
\begin{enumerate}[label=(\alph*), leftmargin=*]
  \item At the assembly level, `counter = counter + 1` is compiled into three separate physical instructions:
        1. **Read:** Load value of `counter` from RAM into CPU register $R_0$.
        2. **Modify:** Increment $R_0$ by 1 ($R_0 := R_0 + 1$).
        3. **Write:** Store value of $R_0$ back to `counter` in RAM.
        If Thread 1 and Thread 2 execute these instructions concurrently, draw a diagram showing an interleaving sequence where both threads read `counter = 0`, modify it, and write back, resulting in a final value of `1` instead of `2`.
  \item This behavior is a **Race Condition**.
        To prevent this, we must define the code block as a **Critical Section** and enforce **Mutual Exclusion**: only one thread can enter the critical section at a time.
        Why does this requirement of mutual exclusion slow down execution on multi-core systems? (Think about thread waiting times).
  \item Modern CPUs provide atomic hardware instructions, such as **Compare-And-Swap (CAS)**:
        `CAS(address, expected_value, new_value)`
        Explain how CAS performs the read, modify, and write steps atomically in a single clock cycle, bypassing race conditions without locking the thread.
\end{enumerate}
\end{puzzlebox}

\subsection{Mutual Exclusion: Semaphores and Mutex Locks}

We formalize coordination primitives and deadlock structures.

\begin{namedconcept}{Dijkstra's Semaphore}
A \newterm{Semaphore} is an integer variable $S$ accessed through two atomic operations:
* $P(S)$ (or wait/down):
  \[
    \text{while } S \le 0 \text{ do } \{ \text{busy wait} \}; \quad S := S - 1
  \]
* $V(S)$ (or signal/up):
  \[
    S := S + 1
  \]
A **Mutex** is a binary semaphore ($S \in \{0, 1\}$) with ownership constraints.
\end{namedconcept}

\begin{explobox}
When threads acquire multiple locks, they can trigger a **Deadlock**, where execution gets permanently stuck.
* **The Coffman Conditions:** A deadlock occurs if and only if four conditions hold simultaneously:
  1. **Mutual Exclusion:** At least one resource is held in a non-shareable mode.
  2. **Hold and Wait:** A thread holds a resource and waits for another.
  3. **No Preemption:** Resources cannot be forcibly taken from threads.
  4. **Circular Wait:** A closed chain of threads exists, each waiting for a resource held by the next.
\begin{enumerate}[label=(\alph*), leftmargin=*]
  \item Prove that if we enforce a global ordering on all locks (e.g. locks must always be acquired in increasing index order: $L_1$ before $L_2$), a deadlock can never occur.
        *(Which of the Coffman conditions is broken by this ordering protocol?)*
  \item Draw a resource allocation graph showing a deadlock state between two threads and two locks.
\end{enumerate}
\end{explobox}

\subsection{Lock-Free structures and the ABA Problem}

We study concurrent data structures without locks.

\begin{explobox}
Locks force threads to sleep, incurring high operating system context-switch latency.
Lock-free algorithms use Compare-And-Swap (CAS) to update nodes dynamically.
* **Treiber Stack Push:**
  \begin{verbatim}
    void push(Node* new_node) {
        do {
            new_node->next = head;
        } while (!CAS(&head, new_node->next, new_node));
    }
  \end{verbatim}
\begin{enumerate}[label=(\alph*), leftmargin=*]
  \item Explain how the Treiber Stack loop handles the case where another thread updates `head` between the assignment `new_node->next = head` and the execution of the `CAS` instruction.
  \item Under concurrent execution, lock-free structures suffer from the **ABA Problem**:
        1. Thread 1 reads value $A$ at `head`, and reads `head->next` pointing to $B$.
        2. Thread 2 pops $A$, pops $B$, and then pushes $A$ back.
        3. Thread 1 resumes and executes `CAS(&head, A, B)`.
        Explain why the CAS succeeds, and why this causes memory corruption if node $B$ has already been deallocated.
\end{enumerate}
\end{explobox}

\newpage

% ─── Section: Discovery ──────────────────────────────────────────────────────
\section{Part 3 --- Discovery}

\begin{discoverybox}
\textbf{What have we built?}

\textbf{First.} Race conditions occur when concurrent threads interleave read-modify-write steps on shared memory variables.

\textbf{Second.} Semaphores and Mutexes enforce mutual exclusion by blocking threads, which can lead to deadlocks if the Coffman conditions are met.

\textbf{Third.} Lock-ordering protocols prevent deadlocks by breaking the Circular Wait condition.

\textbf{Fourth.} Lock-free structures use atomic hardware instructions (CAS) to optimize concurrency, requiring care to prevent ABA memory corruptions.
\end{discoverybox}

\newpage

% ─── Section: Naming Concepts ────────────────────────────────────────────────
\section{Part 4 --- What This Is Called}

\begin{namedconcept}{Concurrency parameters}
* **Race Condition:** An execution anomaly where output is dependent on the sequence or timing of concurrent threads.
* **Critical Section:** A code block that accesses shared mutable resources.
* **Mutual Exclusion:** The requirement that only one thread can execute a critical section at any instant.
* **Deadlock:** A state where threads are blocked waiting for resources held by each other.
* **Compare-And-Swap (CAS):** An atomic CPU instruction used to implement lock-free synchronization.
* **ABA Problem:** A synchronization bug where a memory location is read as $A$, changed to $B$, and changed back to $A$, fooling a CAS check.
\end{namedconcept}

\newpage

% ─── Section: Solved Exercises ───────────────────────────────────────────────
\section{Part 5 --- Solved Exercises}

\begin{problembox}
\begin{enumerate}[label=\textbf{E\arabic*.}]
  \item \textbf{Interleaving Execution Race Condition.}
  
        Draw an interleaving execution trace showing how two threads incrementing a shared `counter` (initially 0) can result in a final value of 1.
        
        \textbf{Solution:}
        Let the steps for Thread 1 be $R_1$ (read), $M_1$ (modify), $W_1$ (write).
        Let the steps for Thread 2 be $R_2$, $M_2$, $W_2$.
        
        Let us trace the execution sequence:
        1. **Step 1: $R_1$ executes.**
           Thread 1 reads `counter = 0` from memory. Register $R_{0,\text{Thread1}} = 0$.
        2. **Step 2: $R_2$ executes.**
           Thread 2 reads `counter = 0` from memory. Register $R_{0,\text{Thread2}} = 0$.
        3. **Step 3: $M_1$ executes.**
           Thread 1 increments its register: $R_{0,\text{Thread1}} = 1$.
        4. **Step 4: $M_2$ executes.**
           Thread 2 increments its register: $R_{0,\text{Thread2}} = 1$.
        5. **Step 5: $W_1$ executes.**
           Thread 1 writes $R_{0,\text{Thread1}} = 1$ to memory. `counter` becomes 1.
        6. **Step 6: $W_2$ executes.**
           Thread 2 writes $R_{0,\text{Thread2}} = 1$ to memory. `counter` becomes 1.
        
        The program ends. The final value of `counter` is 1, even though two increments were executed.
        
  \item \textbf{Lock Acquisition Deadlock Trace.}
  
        Prove that the lock acquisition sequence:
        * Thread 1: Acquires $L_1$, then $L_2$.
        * Thread 2: Acquires $L_2$, then $L_1$.
        can result in a deadlock.
        
        \textbf{Solution:}
        Let us trace an interleaving where a deadlock occurs:
        1. **Step 1:** Thread 1 requests and successfully acquires Lock $L_1$.
        2. **Step 2:** Thread 2 requests and successfully acquires Lock $L_2$.
        3. **Step 3:** Thread 1 requests Lock $L_2$.
           Since Thread 2 holds $L_2$, Thread 1 blocks and waits.
        4. **Step 4:** Thread 2 requests Lock $L_1$.
           Since Thread 1 holds $L_1$, Thread 2 blocks and waits.
           
        Both threads are now permanently blocked.
        Let us verify the **Coffman Conditions**:
        1. *Mutual Exclusion:* Locks $L_1$ and $L_2$ can only be held by one thread at a time. (True).
        2. *Hold and Wait:* Thread 1 holds $L_1$ and waits for $L_2$; Thread 2 holds $L_2$ and waits for $L_1$. (True).
        3. *No Preemption:* Thread 1 cannot force Thread 2 to release $L_2$. (True).
        4. *Circular Wait:* Thread 1 waits for Thread 2, which waits for Thread 1. (True).
        
        Since all 4 conditions hold, the system is deadlocked.
\end{enumerate}
\end{problembox}

\newpage

% ─── Section: Problems ───────────────────────────────────────────────────────
\section{Part 6 --- Problems}

\begin{problembox}
\begin{enumerate}[label=\textbf{P\arabic*.}]
  \item \textbf{Dekker's Algorithm for Mutual Exclusion.}
  
        Dekker's algorithm was the first software-only solution to the mutual exclusion problem for two threads:
        \begin{verbatim}
          // Thread 0
          wants_to_enter[0] = true;
          while (wants_to_enter[1]) {
              if (turn == 1) {
                  wants_to_enter[0] = false;
                  while (turn == 1) { /* busy wait */ }
                  wants_to_enter[0] = true;
              }
          }
          // Critical Section
          turn = 1;
          wants_to_enter[0] = false;
        \end{verbatim}
        \begin{enumerate}[label=(\alph*)]
          \item Explain how the flag array `wants_to_enter` and the integer `turn` ensure that Thread 0 and Thread 1 can never execute their critical sections simultaneously.
          \item Explain how the algorithm prevents **Deadlock** (mutual blocking) and **Starvation** (where one thread is blocked forever).
        \end{enumerate}

  \item \textbf{Dining Philosophers Deadlock Prevention.}
  
        Five philosophers sit around a table. Between each pair of philosophers is a single chopstick.
        To eat, a philosopher must acquire both their left and right chopsticks.
        \begin{enumerate}[label=(\alph*)]
          \item If every philosopher picks up their left chopstick first, show how this satisfies all 4 Coffman conditions, resulting in a deadlock.
          \item Propose a lock-ordering protocol that breaks the Circular Wait condition to prevent deadlock.
                *(Draw the resource allocation graph for your solution).*
        \end{enumerate}
\end{enumerate}
\end{problembox}

\newpage

% ─── Section: Deep Problems ──────────────────────────────────────────────────
\section{Part 7 --- Deep Problems}

\begin{deepbox}
\begin{enumerate}[label=\textbf{D\arabic*.}]
  \item \textbf{Lock-Free Atomic Stack (Treiber Stack).}
  
        A Treiber stack is a thread-safe stack implemented without locks using Compare-And-Swap.
        The `pop` operation is:
        \begin{verbatim}
          Node* pop() {
              Node* old_head;
              do {
                  old_head = head;
                  if (old_head == NULL) return NULL;
              } while (!CAS(&head, old_head, old_head->next));
              return old_head;
          }
        \end{verbatim}
        
        \begin{enumerate}[label=(\alph*)]
          \item Explain how the loop protects stack updates if multiple threads attempt to pop concurrently.
          \item We analyze the **ABA Problem** under concurrent execution:
                * Initially, the stack contains nodes $A \to B \to C$ (head points to $A$).
                * Thread 1 reads `old_head = A`, and reads `old_head->next = B`.
                * Thread 1 is preempted before executing the CAS.
                * Thread 2 pops $A$, pops $B$, and then pushes $A$ back. The stack now contains $A \to C$.
                * Thread 1 resumes and executes `CAS(&head, A, B)`.
                Explain why this CAS succeeds. Show that the head of the stack now points to $B$, which is deallocated, causing memory corruption.
          \item How do modern lock-free systems prevent the ABA problem? (Explain the use of tagged pointers or hazard pointers).
        \end{enumerate}

  \item \textbf{Peterson's Algorithm Correctness Proof.}
  
        Peterson's algorithm is a mutual exclusion algorithm for two threads ($i \in \{0, 1\}$ where $j = 1 - i$):
        \begin{verbatim}
          flag[i] = true;
          turn = j;
          while (flag[j] && turn == j) { /* busy wait */ }
          // Critical Section
          flag[i] = false;
        \end{verbatim}
        
        Prove the correctness of Peterson's algorithm:
        \begin{enumerate}[label=(\alph*)]
          \item Prove **Mutual Exclusion**: Show that both threads cannot be in their critical sections simultaneously.
                *(Hint: Prove by contradiction. Assume both are in their critical sections. Analyze the values of flag[0] and flag[1], and show that the variable `turn` must equal both 0 and 1, which is a contradiction).*
          \item Prove **Progress**: Show that if no thread is in the critical section and a thread wants to enter, it will not be blocked.
        \end{enumerate}
\end{enumerate}
\end{deepbox}

\newpage

% ─── Section: Reflection ─────────────────────────────────────────────────────
\section{Part 8 --- Reflection}

\begin{reflectionbox}
\textbf{Conceptual Reflection}

Consider the following questions:
\begin{enumerate}
  \item **On the cost of safety:** Mutual exclusion guarantees data consistency by forcing threads to wait, sacrificing parallel execution speed. Is there a fundamental physical limit on how fast concurrent cores can coordinate without sacrificing data safety?
  \item **On the nature of lock-freedom:** Lock-free structures prevent threads from sleeping, using CAS loops instead. In highly congested systems, these loops consume CPU cycles in busy-waiting. When is a lock-free structure actually slower than a standard lock?
  \item **On the limits of coordination:** We saw that software-only solutions (like Dekker's or Peterson's algorithms) are complex and rely on strict sequential ordering of memory writes. On modern out-of-order processors, why do these algorithms fail unless we insert hardware **Memory Barriers**?
\end{enumerate}
\end{reflectionbox}

\medskip
\noindent\textbf{Looking ahead.}
Module~6.2 introduces **Memory Models, Cache Coherence, and CPU Boundaries**. Having studied how threads coordinate logically, we will study the physical hardware boundary: memory barriers, store buffers, and cache coherence protocols (MESI) that govern multi-core memory systems.

\vspace{1.4cm}
\begin{center}
{\color{JourneyColor}\rule{0.5\linewidth}{0.8pt}}\\[6pt]
{\small\color{JourneyColor}\textit{End of Journey 6, Module 6.1 --- Concurrency, Mutual Exclusion, and Locks}}\\[4pt]
{\small\color{JourneyColor}Continue to: \textbf{Module 6.2 --- Memory Models, Cache Coherence, and CPU Boundaries}}\\[6pt]
{\color{JourneyColor}\rule{0.5\linewidth}{0.8pt}}
\end{center}

\end{document}
